\documentclass[11pt]{article}
\usepackage[T1]{fontenc}
\usepackage{lmodern}
\usepackage[margin=1in]{geometry}
\usepackage{amsmath,amssymb,amsthm}
\usepackage[hidelinks]{hyperref}
\usepackage{microtype}
\newtheorem{theorem}{Theorem}[section]
\newtheorem{proposition}[theorem]{Proposition}
\newcommand{\R}{\mathbb R}
\newcommand{\E}{\mathbf E}
\newcommand{\lean}[2]{\par\smallskip\noindent\textit{Lean reference:}\;
  \texttt{\detokenize{#1}}\par\noindent
  \hspace*{1em}\texttt{\detokenize{#2}}\par\smallskip}
\title{Clock Regularity and Hybrid Reconstruction\protect\\
for Actual Iterated Brownian Motion}
\author{Project Stardust\protect\\\small Private mathematical synthesis}
\date{Revision 1 --- October 25, 2026}
\begin{document}
\maketitle
\begin{abstract}
For reflected-clock iterated Brownian motion, source averaging produces
multitime covariance functions with folded and collision cusps. We describe
their verified membership in the original closed Gaussian clock derivative
graph, and the actual predictable Brownian-clock representation of completed
Fourier source means. Combining that clock channel with an existing source
adjoint reconstructs the entire centered IBM observable through the original
intrinsic joint divergence. A sine example shows exactly why reconstructing
the mean is insufficient. Hidden-clock predictability, intrinsic observability,
and physical-time stochastic calculus remain distinct requirements.
\end{abstract}

\noindent\textbf{Status.} This note summarizes audited Lean results through
checkpoint 1520, mathematical commit
\begin{quote}\small\texttt{c3381683f4d7abad527dee9575f4cf57deea0431}.\end{quote}
It asserts neither external novelty nor a completed intrinsic IBM calculus.
All theorem names below have prefix
\begin{quote}\small\texttt{OperatorDerivativeRewrite.Foundation.}\end{quote}
Original research manuscripts remain unchanged. All artifacts are private.

\section{The actual process and the averaging operation}
Work on the existing independent scenery--clock Gaussian probability space.
Write $Z_t=X(|Y_t|)$, where $X$ is the constructed Brownian scenery and $Y$ the
constructed Brownian clock. For a finite index set $I$, nonnegative observation
times $r_i$, real frequencies $a_{ki}$ and phases $\beta_k$, let
\[
 F=\sum_{k\geq0}w_k\cos\left(\beta_k+\sum_{i\in I}a_{ki}Z_{r_i}\right),
 \qquad \sum_k|w_k|<\infty.
\]
The preexisting Fourier observation is an actual intrinsic IBM $L^2$ class,
identified almost everywhere with this uniformly convergent scalar series.
Let $M=\E_X F$ be its preexisting source mean. This is a clock Gaussian $L^2$
class. Membership in that space alone would not make it an intrinsic function
of the observed IBM path.

For $b\in\R^I$ and $x\in\R^I$, the genuine scenery covariance gives
\[
 Q_b(x)=\sum_{i,j}b_i b_j\min(|x_i|,|x_j|)
       =\left\|\sum_i b_i h_{|x_i|}\right\|^2\geq0,
\]
where $h_s$ is the actual scenery first-chaos direction. Consequently
\begin{equation}\label{eq:mean}
 M(Y)=\sum_k w_k\cos(\beta_k)
            \exp\bigl(-Q_{a_k}((Y_{r_i})_i)/2\bigr)
 \quad\text{almost everywhere}.
\end{equation}
\lean{IBMFoldedSourceCovariance}{gram; variance_expansion}
\lean{IBMFourierClockMeanRegularity}{clockMean_actual}
This is a law-based source average of the actual observable, not a replacement
random variable defined only through a formal operator identity.

\section{Cusps in the original closed Gaussian graph}
Let $D_Y$ denote the original closed Gaussian derivative and let $e_r=D_Y Y_r$.
The existing normalization gives $\|e_r\|=\sqrt r$. Products below are
almost-everywhere statements about genuine $L^2$ representatives.

\begin{theorem}[Absolute value and minimum]
If $(f,g)$ belongs to the original closed graph, then $|f|$ has graph output
$\operatorname{signum}(f)g$, of norm at most $\|g\|$. If also $(h,k)$ belongs
to that graph, then $\min(f,h)$ has output
\[
 \frac{g+k-\operatorname{signum}(f-h)(g-k)}2,
 \qquad \text{of norm at most }\|g\|+\|k\|.
\]
The convention is $\operatorname{signum}(0)=0$; no null-level hypothesis is
imposed on either zero sets or collision sets.
\end{theorem}
\lean{GaussianAbsoluteMinimumGraph}{absolute_graph; minimum_graph}
Bounded sine clips of smooth absolute radii converge strongly in both graph
coordinates. The proof uses the existing bounded chain rule and closedness of
the original operator. It does not introduce a new weak derivative to avoid
the domain question.

Put $v_r=\operatorname{signum}(Y_r)e_r$ and
\[
 H_{rs}=\frac{v_r+v_s-
   \operatorname{signum}(|Y_r|-|Y_s|)(v_r-v_s)}2.
\]
Then $H_{rs}$ is the actual graph output of $\min(|Y_r|,|Y_s|)$, and
$H_{rr}=v_r$ exactly. The output of $Q_b((Y_{r_i})_i)$ is
$V_b=\sum_{i,j}b_i b_jH_{r_i r_j}$. Writing $A(b)=\sum_i|b_i|$, its norm is
at most
\[
 \sum_{i,j}|b_i b_j|(\sqrt{r_i}+\sqrt{r_j})
 \leq 2A(b)^2\sum_i\sqrt{r_i}.
\]
Signed rows, duplicate clocks, zero times and empty observation sets are
included. For $\exp(-Q_b/2)$ the verified graph output is
\[
 -\tfrac12\exp(-|Q_b|/2)\operatorname{signum}(Q_b)V_b.
\]
Nonnegativity removes the absolute value in the \emph{value}. The output
retains the zero-level sign convention; no second-derivative or divergence
domain assertion is inferred from this first graph rule.
\lean{IBMMultitimeClockVarianceGraph}{overlap_graph; variance_graph; damping_graph}

\begin{theorem}[Completed source-mean graph]\label{thm:graph}
Assume the value and curvature summability conditions
\[
 \sum_k|w_k|<\infty,\qquad \sum_k |w_k|A(a_k)^2<\infty.
\]
Then $M\in\operatorname{Dom}(D_Y)$.
Its output is the actual vector series obtained by multiplying each damping
output above by $w_k\cos\beta_k$ and summing. Both scalar and vector series
are norm summable. The output norm is at most
\[
 \left(\sum_k|w_k|A(a_k)^2\right)\sum_i\sqrt{r_i}.
\]
\end{theorem}
\lean{IBMFourierClockGraph}{multitime_clock_graph; multitime_clock_domain}
The bound is sufficient, not claimed sharp. The stronger single-clock bound
proved earlier is preserved. Classical cusp nondifferentiability therefore
does not obstruct this original first-order closed graph.

\section{A genuine predictable clock identity}
Let $J_Y$ be the independently constructed Brownian-clock It\^o integral on
its actual predictable $L^2$ space, and let $P_Y$ denote the existing
predictable projection of Gaussian vector fields.

\begin{theorem}[Clock PRP and Clark--Ocone]
For every Fourier observation above with $\sum_k|w_k|<\infty$, the actual mean
$M$ lies in the raw Brownian-generated $L^2$ space. The unique predictable
coefficient $U=J_Y^*M$ satisfies
\begin{equation}\label{eq:clock}
 J_YU=M-\E_YM.
\end{equation}
If every $r_i\leq T$, $U$ is supported in $[0,T]$, and the original
finite-horizon It\^o integral gives the same identity. Under the additional
curvature summability in Theorem~\ref{thm:graph},
$U=P_YD_YM$. Coefficients obey
$\|U-U'\|\leq\|M-M'\|$ in their actual $L^2$ norms.
\end{theorem}
\lean{IBMFourierClockClarkOcone}{stochastic_reconstruction; coefficient_unique}
\lean{IBMFourierClockClarkOcone}{coefficient_supported; finite_reconstruction}
\lean{IBMFourierClockClarkOcone}{projected_actual_derivative; coefficient_stability}
The representation, uniqueness and support results do not require curvature
summability; that extra hypothesis identifies the coefficient with the full
original derivative. Brownian and finite-past measurability are proved from
the actual covariance series, not from ambient Gaussian measurability.

\section{The actual IBM observable needs both channels}
Let $K_XF$ be the preexisting bounded source-adjoint recovery. Let $\delta_X$
and $\delta_Y$ be the original transported ambient source and clock adjoints.
Lift the predictable clock field to the joint space by making it constant in
the scenery variable, and call this field $C(U)$. Original adjoint membership
and exact source-vacuum identities give
\[
 \delta_XK_XF=F-M,\qquad
 \delta_YC(U)=M-\E_YM,\qquad \E_YM=\E F,
\]
where clock-only variables are lifted to the product space in these formulas.

\begin{theorem}[Original intrinsic hybrid reconstruction]
The pair $(K_XF,C(U))$ belongs to the domain of the original intrinsic joint
divergence $\delta_{\mathrm{IBM}}$, and
\begin{equation}\label{eq:hybrid}
 \delta_{\mathrm{IBM}}(K_XF,C(U))=F-\E F.
\end{equation}
Its adjoint identity holds against every input in the original intrinsic
joint derivative domain. The reconstruction needs only $\sum_k|w_k|<\infty$.
\end{theorem}
\lean{IBMFourierHybridReconstruction}{source_vacuum; clockMean_expectation}
\lean{IBMFourierHybridReconstruction}{partial_graph; clock_value; reconstruction; duality}
The last projection onto the actual IBM observable space is legitimate
because the centered $F$ is already in that space. There is no assumption
that projection preserves derivative domains or predictable fields.

For a concrete test, take one frequency, one nonzero weight and phase
$\pi/2$. The actual observable is $F=-\sin(Z_R)$. Its source mean and its
entire clock coefficient vanish, but for every $R>0$ it is nonconstant and
outside the original joint derivative domain. Identity~\eqref{eq:hybrid}
still applies through the source channel. Thus~\eqref{eq:clock} cannot be
substituted for reconstruction of $F$ itself.
\lean{IBMFourierClockClarkOcone}{coefficient_zero_of_cos_zero}
\lean{IBMActualSineDomainObstruction}{sinePoint_not_constant; sinePoint_not_joint_domain}
The precise Fourier/sine identification and the simultaneous positive and
negative domain tests are retained in checkpoint 1520's independent Lean checks.

\section{What is still missing}
The clock coefficient is predictable for the \emph{hidden Brownian clock}.
The source channel is the existing static source adjoint, not a claimed
physical-time predictable integrand. Neither construction is automatically
measurable for the observed IBM filtration. Only $U$ has the uniqueness
assertion above; joint divergence integrands can have a kernel.

Other verified results do not remove that distinction. Nonseparable
$C^1$-time/$C^4$-space Simpson formulas and Bernstein continuous-path
realizations hold in their specified probability topologies. The Bernstein
construction uses all grid prefixes and is not asserted causal. Separately,
the original native grid outputs have no weak intrinsic Hilbert limit, even
along a strictly increasing subsequence. Existing Fourier source-lowered
actions can converge strongly without contradicting that obstruction:
they apply a different tested/lowered map with a specified smaller domain.
\lean{IBMTimeSpaceSimpson}{sums_tendstoInMeasure}
\lean{IBMContinuousSimpsonRealization}{approximation_tendstoInMeasure}
\lean{IBMNativeWeakObstruction}{no_original_weak_subsequence}
\lean{IBMFourierCylinderSeries}{native_tendsto; source_native_tendsto}

A precise remaining task is to identify an admissible coefficient in the
\emph{observed} predictable integral space, prove its actual integral agrees
with the desired physical prescription, and show that the centered remainder
is zero in the required topology. Alternatively, an explicit nonzero residual
would delimit that proposed formula. Hidden-clock representation alone
establishes neither conclusion. Nonsmooth physical-time It\^o--Tanaka,
domain/closure transfers and arbitrary-partition claims retain separate proof
obligations. No new campaign, novel-theorem priority or completion claim is
created by this synthesis.
\end{document}
